\documentclass[10pt,a4paper]{amsart}
\usepackage{amsmath,amssymb,amsthm,graphicx}
% For commutative diagrams we use Xy-pic
\usepackage[all]{xy}

% We use 2cell for 2-commutative diagrams.
\xyoption{2cell}
% Theorem environments.
%
\theoremstyle{plain}
\newtheorem{theorem}[subsection]{Theorem}
\newtheorem{proposition}[subsection]{Proposition}
\newtheorem{lemma}[subsection]{Lemma}

\theoremstyle{definition}
\newtheorem{definition}[subsection]{Definition}
\newtheorem{example}[subsection]{Example}
\newtheorem{exercise}[subsection]{Exercise}
\newtheorem{situation}[subsection]{Situation}

\theoremstyle{remark}
\newtheorem{remark}[subsection]{Remark}
\newtheorem{remarks}[subsection]{Remarks}

\numberwithin{equation}{subsection}

% Macros
%
\def\lim{\mathop{\mathrm{lim}}\nolimits}
\def\colim{\mathop{\mathrm{colim}}\nolimits}
\def\Spec{\mathop{\mathrm{Spec}}}
\def\Hom{\mathop{\mathrm{Hom}}\nolimits}
\def\Ext{\mathop{\mathrm{Ext}}\nolimits}
\def\SheafHom{\mathop{\mathcal{H}\!\mathit{om}}\nolimits}
\def\SheafExt{\mathop{\mathcal{E}\!\mathit{xt}}\nolimits}
\def\Sch{\mathit{Sch}}
\def\Mor{\mathop{\mathrm{Mor}}\nolimits}
\def\Ob{\mathop{\mathrm{Ob}}\nolimits}
\def\Sh{\mathop{\mathit{Sh}}\nolimits}
\def\NL{\mathop{N\!L}\nolimits}
\def\CH{\mathop{\mathrm{CH}}\nolimits}
\def\proetale{{pro\text{-}\acute{e}tale}}
\def\etale{{\acute{e}tale}}
\def\QCoh{\mathit{QCoh}}
\def\Ker{\mathop{\mathrm{Ker}}}
\def\Im{\mathop{\mathrm{Im}}}
\def\Coker{\mathop{\mathrm{Coker}}}
\def\Coim{\mathop{\mathrm{Coim}}}

% Boxtimes
%
\DeclareMathSymbol{\boxtimes}{\mathbin}{AMSa}{"02}

%
% Macros for moduli stacks/spaces
%
\def\QCohstack{\mathcal{QC}\!\mathit{oh}}
\def\Cohstack{\mathcal{C}\!\mathit{oh}}
\def\Spacesstack{\mathcal{S}\!\mathit{paces}}
\def\Quotfunctor{\mathrm{Quot}}
\def\Hilbfunctor{\mathrm{Hilb}}
\def\Curvesstack{\mathcal{C}\!\mathit{urves}}
\def\Polarizedstack{\mathcal{P}\!\mathit{olarized}}
\def\Complexesstack{\mathcal{C}\!\mathit{omplexes}}
% \Pic is the operator that assigns to X its picard group, usage \Pic(X)
% \Picardstack_{X/B} denotes the Picard stack of X over B
% \Picardfunctor_{X/B} denotes the Picard functor of X over B
\def\Pic{\mathop{\mathrm{Pic}}\nolimits}
\def\Picardstack{\mathcal{P}\!\mathit{ic}}
\def\Picardfunctor{\mathrm{Pic}}
\def\Deformationcategory{\mathcal{D}\!\mathit{ef}}
\usepackage{cleveref}
\crefname{originallemma}{original lemma}{original lemmas}
\Crefname{originallemma}{Original lemma}{Original lemmas}
\makeatletter
\newlabel{lemma-colimit}{{00D5}{source}{}{sourcepointer-1}{}}
\newlabel{lemma-colimit@cref}{{[originallemma][0][]00D5}{[1][0][]source}{}{}{}}
\makeatother
\crefname{originallemma}{original lemma}{original lemmas}
\Crefname{originallemma}{Original lemma}{Original lemmas}
\makeatletter
\newlabel{lemma-constructible-maps-into-constant-general}{{09Z7}{source}{}{sourcepointer-2}{}}
\newlabel{lemma-constructible-maps-into-constant-general@cref}{{[originallemma][0][]09Z7}{[1][0][]source}{}{}{}}
\makeatother
\crefname{originallemma}{original lemma}{original lemmas}
\Crefname{originallemma}{Original lemma}{Original lemmas}
\makeatletter
\newlabel{lemma-finite-pushforward-commutes-with-base-change}{{0959}{source}{}{sourcepointer-3}{}}
\newlabel{lemma-finite-pushforward-commutes-with-base-change@cref}{{[originallemma][0][]0959}{[1][0][]source}{}{}{}}
\makeatother
\crefname{originallemma}{original lemma}{original lemmas}
\Crefname{originallemma}{Original lemma}{Original lemmas}
\makeatletter
\newlabel{lemma-finite-pushforward-constructible}{{095R}{source}{}{sourcepointer-4}{}}
\newlabel{lemma-finite-pushforward-constructible@cref}{{[originallemma][0][]095R}{[1][0][]source}{}{}{}}
\makeatother
\crefname{originallemma}{original lemma}{original lemmas}
\Crefname{originallemma}{Original lemma}{Original lemmas}
\makeatletter
\newlabel{lemma-proper-over-henselian-and-h1}{{0A5G}{source}{}{sourcepointer-5}{}}
\newlabel{lemma-proper-over-henselian-and-h1@cref}{{[originallemma][0][]0A5G}{[1][0][]source}{}{}{}}
\makeatother
\crefname{originallemma}{original lemma}{original lemmas}
\Crefname{originallemma}{Original lemma}{Original lemmas}
\makeatletter
\newlabel{lemma-torsion-colimit-constructible}{{03SA}{source}{}{sourcepointer-6}{}}
\newlabel{lemma-torsion-colimit-constructible@cref}{{[originallemma][0][]03SA}{[1][0][]source}{}{}{}}
\makeatother
\crefname{originallemma}{original lemma}{original lemmas}
\Crefname{originallemma}{Original lemma}{Original lemmas}
\makeatletter
\newlabel{lemma-vanishing-restriction-injective}{{0A51}{source}{}{sourcepointer-7}{}}
\newlabel{lemma-vanishing-restriction-injective@cref}{{[originallemma][0][]0A51}{[1][0][]source}{}{}{}}
\makeatother
\crefname{originalproposition}{original proposition}{original propositions}
\Crefname{originalproposition}{Original proposition}{Original propositions}
\makeatletter
\newlabel{proposition-finite-higher-direct-image-zero}{{03QP}{source}{}{sourcepointer-8}{}}
\newlabel{proposition-finite-higher-direct-image-zero@cref}{{[originalproposition][0][]03QP}{[1][0][]source}{}{}{}}
\makeatother
\crefname{originalproposition}{original proposition}{original propositions}
\Crefname{originalproposition}{Original proposition}{Original propositions}
\makeatletter
\newlabel{proposition-leray}{{03QC}{source}{}{sourcepointer-9}{}}
\newlabel{proposition-leray@cref}{{[originalproposition][0][]03QC}{[1][0][]source}{}{}{}}
\makeatother
\title{First cohomology of an injective sheaf on the closed fibre of the projective line over a henselian local ring vanishes}
\author{The Stacks Project Authors}
\date{stacks-project:etale-cohomology.tex:local-snapshot-20261011; GNU FDL 1.2 or later}
\begin{document}
\maketitle
\noindent\textit{Editorial note.} All delivered mathematical text is the author's own, copied from the pinned Stacks Project chapter file etale-cohomology.tex; no statement or proof step was rewritten, reordered or removed.\par\medskip
\noindent\textit{Editorial note.} The counted proof is the author's complete argument: the proof is written out in full: the two-affine reduction for q > 1, the constructible approximation of the class, the enlargement and splitting of F, the reduction to a finite pushforward, the Leray and base-change computation on the closed fibre, the lift along the proper-over-henselian comparison, and the final commutative square.\par\medskip
\noindent\textit{Editorial note.} The label lemma-efface-cohomology-on-fibre-by-finite-cover names the counted statement, and the proof environment immediately following it in the source is delivered verbatim as its proof.\par\medskip
\noindent\textit{Editorial note.} The statement and proof are copied from the pinned chapter file; the nine cited inputs are kept as external pointers with their canonical Stacks tags, and the author's use of each is retained where it occurs.\par\medskip
\noindent\textit{Editorial note.} The preamble is the chapter's own native preamble excerpt already used for this source: the Xy-pic setup (native preamble lines 14-18) and the native theorem declarations with the mathematical macros (native preamble lines 154-217).\par\medskip
\noindent\textit{Allowed standard prerequisites (editorial).}\begin{itemize}
\item Injective sheaves of Z/ellZ-modules on X\_etale and their restriction to the closed fibre. Same chapter, injective objects of the category of sheaves of Z/ellZ-modules on the etale site and restriction to X\_0. Supplies the coefficient sheaf I whose cohomology on the closed fibre is shown to vanish, and the injectivity used to extend maps.
\item Gabber's affine vanishing for the restriction to the closed fibre. Same chapter, the vanishing lemma for the restriction of an injective sheaf along a closed immersion into a scheme coverable by two affine opens. Supplies the case q > 1 by the fact that P\textasciicircum{}1\_A is separated and covered by two affine opens.
\item Leray spectral sequence and vanishing of higher direct images for a finite morphism. Same chapter, the Leray spectral sequence for a finite morphism and the vanishing of its higher direct images. Supplies the identification H\textasciicircum{}1(X\_0, f\_*M|\_\{X\_0\}) = H\textasciicircum{}1(Y\_0, M) for the finite morphism f used in the reduction.
\item Comparison of H\textasciicircum{}1 over a proper scheme over a henselian local ring with H\textasciicircum{}1 on the closed fibre. Same chapter, the proper-over-henselian statement identifying the two first cohomology groups. Supplies the identification H\textasciicircum{}1(Y\_0, M) = H\textasciicircum{}1(Y, M) for the proper Y -> Spec(A).
\end{itemize}\medskip
\ifdefined\hypertarget\hypertarget{sourcepointer-1}{}\fi
\noindent\textit{Source reference.} Original lemma 00D5: sources/stacks-etale-cohomology/pinned-original-source.tex, source lines 7427--7427. Prerequisite pointer invoked at proof line(s) 16182; it names a standard fact of the same chapter that the counted argument uses, and every step of the counted proof built on it is retained verbatim in the delivered body.\par\medskip
\ifdefined\hypertarget\hypertarget{sourcepointer-2}{}\fi
\noindent\textit{Source reference.} Original lemma 09Z7: sources/stacks-etale-cohomology/pinned-original-source.tex, source lines 13028--13028. Prerequisite pointer invoked at proof line(s) 16198; it names a standard fact of the same chapter that the counted argument uses, and every step of the counted proof built on it is retained verbatim in the delivered body.\par\medskip
\ifdefined\hypertarget\hypertarget{sourcepointer-3}{}\fi
\noindent\textit{Source reference.} Original lemma 0959: sources/stacks-etale-cohomology/pinned-original-source.tex, source lines 8292--8292. Prerequisite pointer invoked at proof line(s) 16207; it names a standard fact of the same chapter that the counted argument uses, and every step of the counted proof built on it is retained verbatim in the delivered body.\par\medskip
\ifdefined\hypertarget\hypertarget{sourcepointer-4}{}\fi
\noindent\textit{Source reference.} Original lemma 095R: sources/stacks-etale-cohomology/pinned-original-source.tex, source lines 12528--12528. Prerequisite pointer invoked at proof line(s) 16204; it names a standard fact of the same chapter that the counted argument uses, and every step of the counted proof built on it is retained verbatim in the delivered body.\par\medskip
\ifdefined\hypertarget\hypertarget{sourcepointer-5}{}\fi
\noindent\textit{Source reference.} Original lemma 0A5G: sources/stacks-etale-cohomology/pinned-original-source.tex, source lines 16134--16134. Prerequisite pointer invoked at proof line(s) 16219; it names a standard fact of the same chapter that the counted argument uses, and every step of the counted proof built on it is retained verbatim in the delivered body.\par\medskip
\ifdefined\hypertarget\hypertarget{sourcepointer-6}{}\fi
\noindent\textit{Source reference.} Original lemma 03SA: sources/stacks-etale-cohomology/pinned-original-source.tex, source lines 12216--12216. Prerequisite pointer invoked at proof line(s) 16181; it names a standard fact of the same chapter that the counted argument uses, and every step of the counted proof built on it is retained verbatim in the delivered body.\par\medskip
\ifdefined\hypertarget\hypertarget{sourcepointer-7}{}\fi
\noindent\textit{Source reference.} Original lemma 0A51: sources/stacks-etale-cohomology/pinned-original-source.tex, source lines 15259--15259. Prerequisite pointer invoked at proof line(s) 16177; it names a standard fact of the same chapter that the counted argument uses, and every step of the counted proof built on it is retained verbatim in the delivered body.\par\medskip
\ifdefined\hypertarget\hypertarget{sourcepointer-8}{}\fi
\noindent\textit{Source reference.} Original proposition 03QP: sources/stacks-etale-cohomology/pinned-original-source.tex, source lines 8245--8245. Prerequisite pointer invoked at proof line(s) 16213; it names a standard fact of the same chapter that the counted argument uses, and every step of the counted proof built on it is retained verbatim in the delivered body.\par\medskip
\ifdefined\hypertarget\hypertarget{sourcepointer-9}{}\fi
\noindent\textit{Source reference.} Original proposition 03QC: sources/stacks-etale-cohomology/pinned-original-source.tex, source lines 8176--8176. Prerequisite pointer invoked at proof line(s) 16211; it names a standard fact of the same chapter that the counted argument uses, and every step of the counted proof built on it is retained verbatim in the delivered body.\par\medskip
\begin{lemma}
\label{lemma-efface-cohomology-on-fibre-by-finite-cover}
Let $A$ be a henselian local ring. Let $X = \mathbf{P}^1_A$.
Let $X_0 \subset X$ be the closed fibre. Let $\ell$ be a prime
number. Let $\mathcal{I}$ be an injective sheaf of
$\mathbf{Z}/\ell\mathbf{Z}$-modules on $X_\etale$. Then
$H^q_\etale(X_0, \mathcal{I}|_{X_0}) = 0$ for $q > 0$.
\end{lemma}

\begin{proof}
Observe that $X$ is a separated scheme which can be covered by $2$
affine opens. Hence for $q > 1$ this follows from Gabber's affine
variant of the proper base change theorem, see
Lemma \ref{lemma-vanishing-restriction-injective}.
Thus we may assume $q = 1$. Let
$\xi \in H^1_\etale(X_0, \mathcal{I}|_{X_0})$.
Goal: show that $\xi$ is $0$.
By Lemmas \ref{lemma-torsion-colimit-constructible} and
\ref{lemma-colimit} we can find a map $\mathcal{F} \to \mathcal{I}$
with $\mathcal{F}$ a constructible sheaf of
$\mathbf{Z}/\ell\mathbf{Z}$-modules
and $\xi$ coming from an element $\zeta$ of
$H^1_\etale(X_0, \mathcal{F}|_{X_0})$. Suppose we have an injective map
$\mathcal{F} \to \mathcal{F}'$ of sheaves of
$\mathbf{Z}/\ell\mathbf{Z}$-modules on $X_\etale$.
Since $\mathcal{I}$ is injective we can extend the given map
$\mathcal{F} \to \mathcal{I}$ to a map $\mathcal{F}' \to \mathcal{I}$.
In this situation we may replace $\mathcal{F}$ by $\mathcal{F}'$
and $\zeta$ by the image of $\zeta$ in $H^1_\etale(X_0, \mathcal{F}'|_{X_0})$.
Also, if $\mathcal{F} = \mathcal{F}_1 \oplus \mathcal{F}_2$ is a direct sum,
then we may replace $\mathcal{F}$ by $\mathcal{F}_i$
and $\zeta$ by the image of $\zeta$ in $H^1_\etale(X_0, \mathcal{F}_i|_{X_0})$.

\medskip\noindent
By Lemma \ref{lemma-constructible-maps-into-constant-general}
and the remarks above we may assume $\mathcal{F}$
is of the form $f_*\underline{M}$ where $M$ is a finite
$\mathbf{Z}/\ell\mathbf{Z}$-module
and $f : Y \to X$ is a finite morphism of finite presentation
(such sheaves are still constructible by
Lemma \ref{lemma-finite-pushforward-constructible}
but we won't need this).
Since formation of $f_*$ commutes with any base change
(Lemma \ref{lemma-finite-pushforward-commutes-with-base-change})
we see that the restriction of $f_*\underline{M}$ to $X_0$ is
equal to the pushforward of $\underline{M}$ via the induced morphism
$Y_0 \to X_0$ of special fibres. By the Leray spectral sequence
(Proposition \ref{proposition-leray})
and vanishing of higher direct images
(Proposition \ref{proposition-finite-higher-direct-image-zero}),
we find
$$
H^1_\etale(X_0, f_*\underline{M}|_{X_0}) = H^1_\etale(Y_0, \underline{M}).
$$
Since $Y \to \Spec(A)$ is proper we can use
Lemma \ref{lemma-proper-over-henselian-and-h1} to see that
the $H^1_\etale(Y_0, \underline{M})$ is equal to
$H^1_\etale(Y, \underline{M})$. Thus we see that our cohomology
class $\zeta$ lifts to a cohomology class
$$
\tilde\zeta \in H^1_\etale(Y, \underline{M}) = H^1_\etale(X, f_*\underline{M})
$$
However, $\tilde \zeta$ maps to zero in
$H^1_\etale(X, \mathcal{I})$ as $\mathcal{I}$ is injective
and by commutativity of
$$
\xymatrix{
H^1_\etale(X, f_*\underline{M}) \ar[r] \ar[d] &
H^1_\etale(X, \mathcal{I}) \ar[d] \\
H^1_\etale(X_0, (f_*\underline{M})|_{X_0}) \ar[r] &
H^1_\etale(X_0, \mathcal{I}|_{X_0})
}
$$
we conclude that the image $\xi$ of $\zeta$ is zero as well.
\end{proof}

\end{document}
